\documentclass[12pt, a4paper]{article}

% Packages for formatting and styling
\usepackage[utf8]{inputenc}
\usepackage[margin=0.75in]{geometry}
\usepackage{fancyhdr}
\usepackage{graphicx}
\usepackage[table]{xcolor}
\usepackage{listings}
\usepackage{titlesec}
\usepackage{booktabs}
\usepackage{array}
\usepackage{tabularx}
\usepackage{longtable}
\usepackage{multirow}
\usepackage{amsmath}
\usepackage{subcaption}
\usepackage{float}
\usepackage[section]{placeins}
\usepackage[shortlabels]{enumitem}
\usepackage[most]{tcolorbox}
\usepackage{siunitx}
\usepackage{microtype}
\usepackage[htt]{hyphenat} % lets long \texttt names break across lines
\usepackage{hyperref}
\usepackage[capitalise,noabbrev]{cleveref}

% Define custom colors for the document and code blocks
\definecolor{accentcolor}{rgb}{0.0, 0.2, 0.4} % Dark blue for headings and accents
\definecolor{codegreen}{rgb}{0.13, 0.55, 0.13}
\definecolor{codegray}{rgb}{0.5, 0.5, 0.5}
\definecolor{codeorange}{rgb}{0.85, 0.33, 0.1}
\definecolor{backcolour}{rgb}{0.97, 0.97, 0.97}

% Setup hyperlink formatting for TOC and references
\hypersetup{
    colorlinks=true,
    linkcolor=accentcolor,
    filecolor=magenta,
    urlcolor=cyan,
    pdftitle={Assignment Report},
}

% Header and Footer setup
\pagestyle{fancy}
\fancyhf{}
\fancyhead[L]{\textcolor{codegray}{\nouppercase{\leftmark}}} % Section name on the left
\fancyhead[R]{\textcolor{codegray}{CS683}}       % Course name on the right
\fancyfoot[C]{\textbf{\thepage}}                              % Page number at the bottom center
\renewcommand{\headrulewidth}{0.4pt}
\renewcommand{\footrulewidth}{0.4pt}
\newtcolorbox{aestheticquote}[1][]{%
    enhanced,
    breakable,
    frame hidden,
    boxrule=0pt,
    colback=backcolour, % Uses the light gray defined in the template
    borderline west={3pt}{0pt}{accentcolor}, % Adds a thick left border using the theme's dark blue
    sharp corners,
    left=12pt,
    right=12pt,
    top=10pt,
    bottom=10pt,
    fontupper=\itshape, % Automatically italicizes the quote text
    #1
}
% Section formatting for a professional look
\titleformat{\section}
  {\Large\bfseries\color{accentcolor}}
  {\thesection}{1em}{}[\vspace{-0.5ex}\rule{\textwidth}{0.5pt}]
\titleformat{\subsection}
  {\large\bfseries\color{accentcolor!80}}
  {\thesubsection}{1em}{}

% Python Code Listing setup
\lstdefinestyle{pythonstyle}{
    language=Python,
    backgroundcolor=\color{backcolour},
    commentstyle=\color{codegreen}\itshape,
    keywordstyle=\color{accentcolor}\bfseries,
    numberstyle=\tiny\color{codegray},
    stringstyle=\color{codeorange},
    basicstyle=\ttfamily\small,
    breakatwhitespace=false,
    breaklines=true,
    captionpos=b,
    keepspaces=true,
    numbers=left,
    numbersep=10pt,
    showspaces=false,
    showstringspaces=false,
    showtabs=false,
    tabsize=4,
    frame=L,           % Left border line
    framerule=2pt,
    rulecolor=\color{accentcolor} % Matches the theme
}
\lstset{style=pythonstyle}
% C++ Code Listing setup
% Modern C++ Code Environment using tcolorbox
\newtcblisting{cppcode}[1][]{
    listing only,
    listing options={
        language=C++,
        basicstyle=\ttfamily\small,
        keywordstyle=\color{accentcolor}\bfseries,
        commentstyle=\color{codegreen}\itshape,
        stringstyle=\color{codeorange},
        numberstyle=\tiny\color{codegray},
        numbers=left,
        numbersep=12pt,
        showstringspaces=false,
        breaklines=true,
        tabsize=4,
    },
    enhanced,
    boxrule=0.8pt,                 % Thin border
    colback=backcolour,            % Background color
    colframe=accentcolor!80,       % Border color
    drop fuzzy shadow,             % Aesthetic drop shadow
    arc=6pt,                       % Rounded corners
    left=5pt, right=5pt, top=2pt, bottom=2pt,
    title={#1},                    % Dynamic title
    fonttitle=\bfseries\sffamily\footnotesize,
    coltitle=white,
    attach boxed title to top left={xshift=15pt, yshift=-3mm},
    boxed title style={size=small, colback=accentcolor, arc=3pt, boxrule=0pt}
}

% ---- small helpers used by the tables and figures of this report ----
% plain booktabs tables: no zebra rows, no tinted header
\colorlet{tablehead}{white}
\colorlet{muted}{codegray}
\newcommand{\tablesetup}{\small\renewcommand{\arraystretch}{1.15}\setlength{\tabcolsep}{4pt}}
\newcommand{\colhead}[1]{\textbf{#1}}
\newcolumntype{L}{>{\raggedright\arraybackslash}X}
\newcommand{\tbd}{--}
\newcommand{\todo}[1]{\textcolor{codegray}{\textit{[#1]}}}
% \resultfig[width]{file-stem}{caption}{label}: figures/<file-stem>.png (or .pdf),
% or a grey box if the plot has not been generated yet
\newcommand{\resultfig}[4][0.75\linewidth]{%
  \begin{figure}[htbp]\centering
  \IfFileExists{figures/#2.pdf}{\includegraphics[width=#1]{figures/#2.pdf}}{%
  \IfFileExists{figures/#2.png}{\includegraphics[width=#1]{figures/#2.png}}{%
  \fbox{\parbox[c][4cm][c]{#1}{\centering\small Plot: \texttt{figures/\detokenize{#2}.png}}}}}%
  \caption{#3}\label{#4}
  \end{figure}}

\setlength{\emergencystretch}{3em} % avoids lines running into the margin with long \texttt paths

\begin{document}

% -------------------------------------------------------------------
% FRONT PAGE
% -------------------------------------------------------------------
\begin{titlepage}
    \centering
    \vspace*{1.5cm}

    {\huge\bfseries\color{accentcolor} Programming Assignment 2\par}
    \vspace{0.5cm}
    {\Large Data Prefetching and Cache Replacement in ChampSim\par}

    \vspace{2cm}

    \vspace{1.5cm}

    {\Large\bfseries Team VS CoPiLot \par}
    \vspace{0.5cm}
    {\large Advay Bapat 23B1225\\Aniket Gupta 23B1245\\ Apoorv Goyal 23B1266\\ Kundan Maiti 23B1242}\\
    {\large Department of Electrical Engineering \par}
    {\large Indian Institute of Technology Bombay \par}

    \vfill

    {\large \textbf{CS683: Advanced Computer Architecture} \par}
    \vspace{0.3cm}

    \vspace{0.3cm}
    {\large \textbf{Date:} \today \par}

    \vspace{1.5cm}
\end{titlepage}

\newpage

% -------------------------------------------------------------------
% TABLE OF CONTENTS
% -------------------------------------------------------------------
% Setting link color to black just for the TOC for a cleaner look
{\hypersetup{linkcolor=black}
\tableofcontents
}
\newpage


% ==========================================
% OVERVIEW
% ==========================================
\section{Overview and Methodology}\label{sec:overview}

Our Stride Pattern
Prefetcher (SPP) keeps, for every load IP, a \emph{short} (one-stride)
and a \emph{long} (three-stride signature) predictor, chooses between them per
IP with a tournament counter, walks the predicted stride path up to eight
predicted lines ahead inside the 4\,KB page, and throttles its degree from measured accuracy.
The final design (\texttt{spp16}, \SI{7.8}{\kilo\byte} of state) reaches a
geometric-mean speedup of \textbf{1.0129} over the no-prefetch baseline,
\textbf{5.4$\times$ the gain of IP-stride} (1.0024).

\subsection{Simulation Setup}\label{sec:setup}
All experiments use the ChampSim simulator from the PA2 repository, with
\num{25}M warm-up instructions followed by \num{25}M simulated instructions
on each of the three provided traces (\texttt{Trace1}, \texttt{Trace2},
\texttt{Trace3}). \Cref{tab:config} lists the simulated system
(values from \texttt{inc/cache.h}).

\begin{table}[htbp]
\centering
\caption{Simulated single-core system.}\label{tab:config}
\tablesetup
\begin{tabular}{ll}
\toprule
\rowcolor{tablehead}\colhead{Component} & \colhead{Configuration} \\
\midrule
Branch predictor & Hashed perceptron \\
L1I cache & 32\,KB, 64 sets $\times$ 8 ways, 4-cycle latency \\
L1D cache & 48\,KB, 64 sets $\times$ 12 ways, 5-cycle latency, 16 MSHRs, 16-entry PQ \\
L2 cache  & 512\,KB, 1024 sets $\times$ 8 ways, 10-cycle latency, 32 MSHRs \\
LLC       & 2\,MB, 2048 sets $\times$ 16 ways, 20-cycle latency \\
DRAM      & 4\,GB, 1 channel, 64-bit, 6400\,MT/s \\
Block / page size & 64\,B / 4\,KB \\
Warm-up / simulation & 25M / 25M instructions \\
\bottomrule
\end{tabular}
\end{table}

% \subsection{Configurations Evaluated}\label{sec:configs}
% \begin{table}[htbp]
% \centering
% \caption{Binaries built for this report (\texttt{./build\_champsim.sh <L1D prefetcher> <L2C replacement> <name>}).}\label{tab:binaries}
% \tablesetup
% \begin{tabular}{llll}
% \toprule
% \rowcolor{tablehead}\colhead{Name} & \colhead{L1D prefetcher} & \colhead{L2C replacement} & \colhead{Used in} \\
% \midrule
% Baseline        & \texttt{no}        & \texttt{lru}       & All tasks \\
% IP-stride       & \texttt{ip\_stride} & \texttt{lru}      & Task 1 \\
% SPP             & \texttt{spp}       & \texttt{lru}       & Tasks 1, 3 \\
% RRIP            & \texttt{no}        & \texttt{rrip}      & Task 2 \\
% I-RRIP          & \texttt{no}        & \texttt{irrip}     & Task 2A \\
% Dyn.\ I-RRIP    & \texttt{no}        & \texttt{dyn\_irrip} & Tasks 2B, 3 \\
% SPP + Dyn.\ I-RRIP & \texttt{spp}    & \texttt{dyn\_irrip} & Task 3 \\
% \bottomrule
% \end{tabular}
% \end{table}

\subsection{Metrics}\label{sec:metrics}
\begin{align}
\text{Speedup} &= \frac{\text{IPC with optimisation}}{\text{IPC of baseline (no prefetching)}} \label{eq:speedup}\\[4pt]
\text{L1D Coverage} &= \frac{\text{L1D load misses}_{\text{baseline}} - \text{L1D load misses}_{\text{prefetcher}}}{\text{L1D load misses}_{\text{baseline}}} \label{eq:coverage}\\[4pt]
\text{Prefetch Accuracy} &= \frac{\text{useful prefetches}}{\text{prefetches issued to the lower level}} \label{eq:accuracy}\\[4pt]
\text{MPKI} &= \frac{\text{load misses}}{\text{instructions}/1000} \label{eq:mpki}
\end{align}
MPKI and prefetch accuracy are taken directly from the ChampSim statistics. Averages across traces are
geometric means for speedup.

\paragraph{Workflow.} Every version was built and run with
\texttt{artifacts/run\_spp.py}, which snapshots the prefetcher source to
\texttt{SPP\_prefetchers/<name>.l1d\_pref}, builds it, logs a note in
\texttt{runsTrackers.txt} and stores the outputs in
\texttt{results\_pf<trace>/<name>.txt}. \texttt{artifacts/parse\_results.py}
turns all outputs into \texttt{results\_summary.csv}, and
\texttt{artifacts/plot\_results.py} draws the figures of this report.

% ==========================================
% TASK 1
% ==========================================
\section{Task 1: Data Prefetching}\label{sec:task1}

\subsection{Task 1A: Understanding the IP-Stride Prefetcher}\label{sec:task1a}
\texttt{prefetcher/ip\_stride.l1d\_pref} is invoked by the L1D on every demand
\emph{load} (hit or miss; stores and prefetches do not train it) with the
virtual address and the IP of the load.

\paragraph{Tracker table.} A fully associative table of 24 \texttt{IP\_TRACKER}
entries, each holding the IP tag, the last cache-line address, the last stride
and an LRU counter. Only the low 8 bits of the IP are kept
(\texttt{MASK\_IP = 0xFF}), so loads whose IPs differ by a multiple of 256
alias to the same entry; the cheap tag is traded for occasional pollution. On a
miss in the table the LRU entry is replaced, its address recorded, and nothing
is prefetched.

\paragraph{Training.} Addresses are converted to cache-line granularity
(\texttt{addr >> 6}), so accesses within the same 64\,B line give a stride of 0
and are ignored. The signed stride to the previous line of the same IP is
computed with an explicit sign (the variables are unsigned), then the entry's
last address and last stride are updated and the entry becomes MRU.

\paragraph{Prefetch issue.} A stride is trusted only when it equals the
previous stride of that IP, i.e.\ prefetching starts on the third access of a
constant-stride stream; there is no confidence counter, so a single irregular
access costs two prefetch opportunities. The prefetcher then issues
\texttt{PREFETCH\_DEGREE = 3} lines at $+s, +2s, +3s$, stopping at the first
one outside the current 4\,KB page (virtual pages are not contiguous in
physical memory). Lines are filled into the L1D while the L1D MSHR is less
than half full and only into the L2 otherwise, so prefetches never starve
demand misses of MSHR entries. \texttt{prefetch\_line} drops the request
silently when the 16-entry prefetch queue is full.

\paragraph{Limitations.} One stride per IP cannot describe a repeating
sequence such as $(+1,\allowbreak+2,\allowbreak+2,\allowbreak+2,\allowbreak+3)^n$: the stride check fails at every change.
It also confirms on two equal strides only, holds just 24 IPs, and prefetches
at most three lines ahead, which is late for streams with long memory latency.
These limitations motivate the design below.

\subsection{Task 1B: Stride Pattern Prefetcher (SPP)}\label{sec:task1b}

\subsubsection{Final Design (\texttt{spp16})}\label{sec:spp-design}
\Cref{tab:spp-structs} lists the state. Every tracked IP owns two small pattern
tables that map a history of recent strides to the stride that followed
it: a short table whose history is the last stride ($K_s=1$) and a long table
whose history (the \emph{signature}) is the last three strides ($K_\ell=3$).
A constant stride $s$ is simply the short-table entry $s\to s$, so the short
table covers uniform strides (and two-stride alternations), while the long
table separates non-uniform sequences such as $(+1,+2,+2,+2,+3)^n$, in which
the stride after $+2$ depends on how many $+2$s came before.

\begin{table}[htbp]
\centering
\caption{State of the final SPP (\texttt{spp16}).}\label{tab:spp-structs}
\tablesetup
\begin{tabularx}{\linewidth}{lrL}
\toprule
\rowcolor{tablehead}\colhead{Structure} & \colhead{Entries} & \colhead{Contents} \\
\midrule
IP tracker & 128 & 8-bit IP tag, last line address, last 3 strides, strides-seen (0--3), 2-bit chooser, LRU \\
Short table (per IP) & 8 & history = 1 stride $\to$ predicted next stride, valid, LRU \\
Long table (per IP) & 8 & signature = 3 strides $\to$ predicted next stride, valid, LRU \\
Throttling & 1 & window counters (fetched, useful), current degree \\
\bottomrule
\end{tabularx}
\end{table}

\paragraph{Per-access algorithm.} On every L1D demand load from a tracked IP
with a non-zero stride $s$:
\begin{enumerate}[nosep]
  \item \textbf{Grade.} Look up the entries matching the \emph{old} history in
        both tables; each predicted this access and is correct if its stored
        stride equals $s$. A missing entry counts as wrong.
  \item \textbf{Choose.} Update the per-IP 2-bit chooser: $+1$ if only the long
        table was right, $-1$ if only the short table was right. New IPs start
        at 0 (short table) so a constant stride is predicted from the third
        access.
  \item \textbf{Train} both tables with (old history $\to s$): overwrite the
        prediction of a matching entry or allocate the LRU entry. The long
        table trains only once three strides have been seen (no zero-padded
        histories).
  \item \textbf{Shift} $s$ into the history.
  \item \textbf{Throttle.} If the current window is complete, recompute the degree.
  \item \textbf{Walk.} Use the long table if chooser $\ge 2$, else the short
        table; if it has no entry for the current history the other table is
        tried (fallback). Starting from the current line, repeatedly add the
        predicted stride, issue a prefetch, and shift the predicted stride into
        a copy of the history, for up to degree steps or until the
        address leaves the 4\,KB page or no entry matches. Fill level follows
        the IP-stride rule (L1D if the MSHR is less than half full, else L2).
\end{enumerate}
The chooser makes the non-uniform mode a reversible, per-IP fall-back: the
uniform (short) predictor has priority, and an IP moves to the signature only
while the signature predicts it better. \Cref{tab:spp-req} maps the
requirements of the assignment to the design.

\begin{table}[htbp]
\centering
\caption{Assignment requirements and how the final SPP meets them.}\label{tab:spp-req}
\tablesetup
\begin{tabularx}{\linewidth}{lL}
\toprule
\rowcolor{tablehead}\colhead{Requirement} & \colhead{Implementation} \\
\midrule
Per-IP pattern table & 128-entry IP tracker, each with its own short and long table \\
Uniform vs.\ non-uniform & short table ($s\to s$) vs.\ 3-stride signature table \\
Uniform first, fall back & chooser starts at short; moves to the signature only when it predicts better, and back \\
Signature $\to$ next stride & long table maps the last 3 strides to the next stride \\
Degree $>1$ & lookahead walk of up to 8 lines along the predicted path \\
Accuracy-based throttling & degree 8/4/2 from windowed accuracy (\cref{sec:spp-throttle}) \\
Page boundary & every candidate is checked against the page of the trigger access \\
\bottomrule
\end{tabularx}
\end{table}

% \begin{cppcode}[Core of \texttt{l1d\_prefetcher\_operate} (\texttt{prefetcher/spp.l1d\_pref}, abridged)]
% // 1. grade both predictors (entries that predicted THIS access)
% bool long_ready = (t.strides_seen >= K_LONG);
% TableEntry* entry_short = matchHistory(t.shortTable, t.lastKstrides);
% TableEntry* entry_long  = long_ready ? matchHistory(t.longTable, t.lastKstrides) : nullptr;
% bool short_correct = entry_short && entry_short->pred_stride == stride;
% bool long_correct  = entry_long  && entry_long->pred_stride  == stride;
% // 2. tournament chooser
% if (long_correct && !short_correct)      { if (t.chooser < CHOOSER_MAX) t.chooser++; }
% else if (short_correct && !long_correct) { if (t.chooser > 0) t.chooser--; }
% // 3. train both tables, 4. shift history
% trainTable(t.shortTable, entry_short, t.lastKstrides, stride);
% if (long_ready) trainTable(t.longTable, entry_long, t.lastKstrides, stride);
% ShiftStride(t.lastKstrides, stride, t.lastKstrides);
% // 5. accuracy-based throttling (sets cur_degree)
% updateThrottle(pf_useful, pf_lower_level);
% // 6. lookahead walk along the predicted path
% bool use_long = (t.chooser >= CHOOSER_THRESHOLD) && (t.strides_seen >= K_LONG);
% Table &chosen = use_long ? t.longTable : t.shortTable;
% Table &other  = use_long ? t.shortTable : t.longTable;
% int64_t pattern[K_LONG];  copy(t.lastKstrides, pattern);
% uint64_t line = cl_addr;
% for (int p = 0; p < cur_degree; ++p) {
%     TableEntry* e = matchHistory(chosen, pattern);
%     if (!e && USE_FALLBACK) e = matchHistory(other, pattern);
%     if (!e) break;
%     uint64_t pf_addr = (line + e->pred_stride) << LOG2_BLOCK_SIZE;
%     line = pf_addr >> LOG2_BLOCK_SIZE;
%     if ((pf_addr >> LOG2_PAGE_SIZE) != (addr >> LOG2_PAGE_SIZE)) break;
%     prefetch_line(ip, addr, pf_addr, MSHR.occupancy < MSHR.SIZE/2 ? FILL_L1 : FILL_L2, 0);
%     ShiftStride(pattern, e->pred_stride, pattern);
% }
% \end{cppcode}

\subsubsection{How the Final Design Was Reached}\label{sec:spp-evolution}
Each version was run on all three traces.
\Cref{tab:evolution} summarises the path; the complete statistics are in
\cref{tab:full}.

\begin{table}[htbp]
\centering
\caption{Design history (speedup over the no-prefetch baseline). Storage for spp14--16 from \cref{sec:spp-storage}.}\label{tab:evolution}
\tablesetup
\footnotesize
\begin{tabularx}{\linewidth}{lLrrrr}
\toprule
\rowcolor{tablehead}\colhead{Version} & \colhead{Change} & \colhead{T1} & \colhead{T2} & \colhead{T3} & \colhead{Geomean} \\
\midrule
ip\_stride & reference (24 IPs, degree 3) & 1.0044 & 1.0007 & 1.0021 & 1.0024 \\
spp1  & per-IP pattern table, $K=3$, degree 1 & 1.0020 & -- & -- & -- \\
spp2  & lookahead walk, degree 4 & 1.0033 & 1.0011 & 1.0029 & 1.0024 \\
spp3  & spp2 with $K=1$ & 1.0052 & 1.0007 & 1.0043 & 1.0034 \\
spp4  & escalating dynamic $K$, 2-bit confidence, off switch & 1.0049 & 1.0007 & 1.0035 & 1.0030 \\
spp5  & spp4, 64 trackers & -- & 1.0049 & -- & -- \\
spp6  & spp4, 256 trackers & 1.0007 & 1.0026 & 1.0001 & 1.0011 \\
spp7  & spp4, 64 trackers, degree 8 & 1.0058 & 1.0060 & 1.0100 & 1.0073 \\
spp8  & spp3, 256 trackers, degree 8 & 1.0070 & 1.0147 & 1.0179 & 1.0132 \\
spp9  & spp8 with $K=3$, 64 entries/IP & 1.0060 & 1.0155 & 1.0095 & 1.0103 \\
spp10 & tournament $K=1$ / $K=3$, fallback on & 1.0073 & 1.0148 & 1.0183 & 1.0134 \\
spp11 & spp10, fallback off & 1.0073 & 1.0148 & 1.0177 & 1.0133 \\
spp12 & + throttling, degrees 16/12/8/4/2 & 1.0065 & 1.0149 & 1.0204 & 1.0139 \\
spp14 & throttling capped at 8 (8/4/2) & 1.0079 & 1.0157 & 1.0196 & \textbf{1.0144} \\
spp15 & spp14, 128 trackers, 8/16 entries (12\,KB) & 1.0081 & 1.0141 & 1.0165 & 1.0129 \\
\textbf{spp16} & \textbf{spp15, long table 8 entries (7.8\,KB)} & \textbf{1.0080} & \textbf{1.0141} & \textbf{1.0165} & \textbf{1.0129} \\
\bottomrule
\end{tabularx}
\end{table}

\paragraph{Step 1 -- lookahead (spp1 $\to$ spp2).} The first version kept, per
IP, a table from the last three strides to the next one and issued one
prefetch. Degree 1 is too late (27\% of useful prefetches late on Trace~1);
walking the predicted path four steps ahead (each step adding the predicted
stride to the previous \emph{predicted} line, not to the trigger) roughly
doubled coverage (2.3\% $\to$ 4.1\%) and halved lateness.

\paragraph{Step 2 -- longer history is not better (spp2 vs.\ spp3).} With only
$K$ changed, $K=1$ beat $K=3$ on coverage on every trace (7.1 vs.\ 4.1\% on
Trace~1, 6.2 vs.\ 3.3\% on Trace~3) while $K=3$ was more accurate (55 vs.\ 48\%,
57 vs.\ 33\% on Trace~2). A three-stride history is a narrower key: it warms up
two strides later, one irregular stride spoils the next three histories, and
the walk stops at the first unseen signature. The predictable misses in these
traces are mostly uniform or short patterns.

\paragraph{Step 3 -- escalating $K$ fails (spp4--spp7, brief).} We made $K$
dynamic per IP: start at 1, raise $K$ when an entry's error counter saturates,
switch the IP off at maximum $K$. It never beat fixed $K=1$, and with 256
trackers (spp6) it collapsed (geomean 1.0011, 2{,}460 lines fetched on Trace~1
instead of $\approx$40{,}000). Out-of-order load reordering and loop restarts
look like conflicts, so IPs escalated, cleared their tables and never came
back down; with 24 trackers the frequent tracker evictions had been resetting
them and hiding the problem. Lesson: any change of history length must be
reversible and must not discard trained state.

\paragraph{Step 4 -- capacity and timeliness (spp8, spp9).} Two findings from
step 3 carried over: more trackers help Trace~2, whose large code footprint
(L1I MPKI 40) means many live load IPs (24 $\to$ 64 trackers: 0.8 $\to$ 1.9\%
coverage), and degree 8 removes most late prefetches (Trace~3: 32\% $\to$ 11\%
late). Applying both to the simple $K=1$ design gave spp8, the first large jump
(geomean 1.0024 $\to$ 1.0132; Trace~3 coverage 17.1\%). The same configuration
with $K=3$ (spp9) was worse overall (1.0103) but \emph{better} on Trace~2
(1.0155 vs.\ 1.0147, accuracy 71 vs.\ 42\%): no single history length is best
for every workload.

\paragraph{Step 5 -- tournament chooser (spp10, spp11).} To get $K=1$ where it
wins and $K=3$ where it wins, both tables are kept and trained, and a per-IP
2-bit counter chooses, graded on every access. \Cref{tab:chooser} shows the
behaviour with the large tables. spp10 reached the best coverage on every trace
(12.2/5.1/17.6\%) and a geomean of 1.0134, never worse than spp8. On a
synthetic $(+1,+2,+2,+2,+3)^n$ IP it prefetched 96.9\% of accesses versus
76.3\% for $K=1$ alone, while matching $K=1$ on uniform streams. It did not
reproduce spp9's precision on Trace~2: there both predictors are right for only
$\approx$11\% of next strides, so spp9's advantage came from \emph{not}
prefetching when its signature was unseen. Turning the fallback off (spp11)
recovered almost none of it (accuracy 45 $\to$ 47\%, geomean 1.0133), so the
fallback was kept and the remaining waste was attacked with throttling.

\begin{table}[htbp]
\centering
\caption{Chooser statistics of spp14 (24/64 entries per IP); counters include warm-up.}\label{tab:chooser}
\tablesetup
\begin{tabular}{lrrr}
\toprule
\rowcolor{tablehead}\colhead{Statistic} & \colhead{Trace 1} & \colhead{Trace 2} & \colhead{Trace 3} \\
\midrule
Walks driven by the long table & 12.1\% & 30.7\% & 33.6\% \\
Next stride predicted correctly: short / long & 22.0 / 18.3\% & 10.8 / 10.9\% & 58.2 / 57.0\% \\
Chooser flips (each direction) & $\approx$15.4k & $\approx$17.6k & $\approx$29.9k \\
\bottomrule
\end{tabular}
\end{table}

\paragraph{Step 6 -- accuracy throttling (spp12 $\to$ spp14).} Detailed in
\cref{sec:spp-throttle}: throttling \emph{down} in inaccurate phases helped
Trace~3 strongly, throttling \emph{up} beyond 8 did not help anywhere, so the
final bands cap the degree at 8 (geomean 1.0144).

\paragraph{Step 7 -- storage (spp14 $\to$ spp16).} spp14 needs $\approx$88\,KB,
more than the 48\,KB L1D. Halving the trackers to 128 and shrinking the per-IP
tables to 8/16 entries (spp15, 12\,KB) cost 0.15 points of geomean, mostly on
Trace~3 (1.0196 $\to$ 1.0165) and Trace~2 (1.0157 $\to$ 1.0141), while
Trace~1 was unchanged. Shrinking the long table further to 8 entries (spp16,
7.8\,KB) changed nothing measurable (geomean 1.0129 in both). With small
tables the long table drives only 2.5/0.7/2.0\% of walks, so most of the
benefit comes from the short predictor, lookahead and throttling. spp16 is the
final design: about 90\% of spp14's gain for 9\% of its storage.

\subsubsection{Accuracy-Based Throttling}\label{sec:spp-throttle}
The prefetcher reads the L1D's own counters \texttt{pf\_useful} and
\texttt{pf\_lower\_level}. A \emph{window} ends every 256 prefetches that
fetched a line from L2; its accuracy is
$\Delta\texttt{pf\_useful}/\Delta\texttt{pf\_lower\_level}$ over the window,
and it sets the degree used by all walks until the next window ends
(\cref{tab:spp-degree}). The first window uses degree 8. ChampSim resets these
counters at the end of warm-up; a counter that goes backwards simply starts a
new window. In hardware the same result needs only two small window counters
(\cref{tab:spp-storage}).

\begin{table}[htbp]
\centering
\caption{Prefetch degree chosen for each window-accuracy range (final design).}\label{tab:spp-degree}
\tablesetup
\begin{tabularx}{\linewidth}{lcL}
\toprule
\rowcolor{tablehead}\colhead{Window accuracy} & \colhead{Degree} & \colhead{Reason} \\
\midrule
$\text{acc} \ge 30\%$          & 8 & distance needed for timeliness; deeper walks only add redundant or page-cut requests \\
$15\% \le \text{acc} < 30\%$  & 4 & poor phase: halve the walk, keep the most accurate first steps \\
$\text{acc} < 15\%$           & 2 & bad phase: keep sampling (degree 0 could never measure recovery) \\
\bottomrule
\end{tabularx}
\end{table}

\paragraph{How the thresholds were chosen.} The data showed that low accuracy
alone is not a reason to throttle: Trace~3 runs at only 22--24\% accuracy at
degree 8, yet going from degree 4 to 8 was its largest single gain (spp4
1.0035 $\to$ spp7 1.0100) because lateness fell from 32\% to 11\%. A threshold
much above 30\% would therefore throttle Trace~3's useful phases. At 30\% the
average window accuracy (50/65/41\% on the three traces) stays above the
threshold, so the common case keeps degree 8, and only clearly worse phases are
cut. Below 15\%, more than 85\% of fetched lines are never used, which is pure
pollution of the 48\,KB L1D, so the degree is reduced again. The window of 256
fetched prefetches gives 470--950 decisions per run; at 50\% accuracy its
sampling error is about $\pm3\%$, small next to the 15-point band width.

\paragraph{Why these degrees.} spp12 used five bands up to degree 16
($\ge$60\% $\to$ 16, 40--60\% $\to$ 12, 20--40\% $\to$ 8, 10--20\% $\to$ 4,
$<$10\% $\to$ 2). \Cref{tab:throttle} compares it with the final bands. On
Trace~2, degree 16 was used in 51\% of windows, raised requests from 2.06M to
2.49M, but fetched and useful lines did not change: the extra steps were
already cached, merged in the prefetch queue, or cut at the page boundary. On
Trace~1, the 57\% of windows at degree 12--16 added 5.7k fetched but only
0.8k useful lines (speedup 1.0073 $\to$ 1.0065). Capping at 8 (spp14) fixed
Trace~1 (1.0079), improved Trace~2 (1.0157) and kept almost all of Trace~3's
gain (1.0196 vs.\ 1.0204), with the highest accuracy of all versions. Degrees
4 and 2 halve the walk at each lower band, which keeps the most accurate first
steps while still producing fetches from which the next window can measure
recovery.

\begin{table}[htbp]
\centering
\caption{Effect of the throttling bands (share of windows at each degree; speedup; accuracy).}\label{tab:throttle}
\tablesetup
\footnotesize
\begin{tabular}{llrrrrrrr}
\toprule
\rowcolor{tablehead}\colhead{Version} & \colhead{Trace} & \colhead{deg 2} & \colhead{deg 4} & \colhead{deg 8} & \colhead{deg 12} & \colhead{deg 16} & \colhead{Speedup} & \colhead{Accuracy} \\
\midrule
spp10 (fixed 8) & 1 & -- & -- & 100\% & -- & -- & 1.0073 & 44.5\% \\
spp12 (to 16)   & 1 & 0.8\% & 4.6\% & 37.7\% & 33.4\% & 23.5\% & 1.0065 & 43.1\% \\
spp14 (to 8)    & 1 & 1.4\% & 13.6\% & 85.0\% & -- & -- & 1.0079 & 46.6\% \\
\midrule
spp10 (fixed 8) & 2 & -- & -- & 100\% & -- & -- & 1.0148 & 45.3\% \\
spp12 (to 16)   & 2 & 1.3\% & 5.4\% & 27.7\% & 14.4\% & 51.2\% & 1.0149 & 45.2\% \\
spp14 (to 8)    & 2 & 1.6\% & 15.6\% & 82.8\% & -- & -- & 1.0157 & 50.2\% \\
\midrule
spp10 (fixed 8) & 3 & -- & -- & 100\% & -- & -- & 1.0183 & 24.0\% \\
spp12 (to 16)   & 3 & 7.5\% & 15.6\% & 39.3\% & 17.8\% & 19.8\% & 1.0204 & 33.5\% \\
spp14 (to 8)    & 3 & 10.5\% & 26.8\% & 62.7\% & -- & -- & 1.0196 & 35.2\% \\
\bottomrule
\end{tabular}
\end{table}

\subsubsection{Storage Overhead}\label{sec:spp-storage}
\Cref{tab:spp-storage} counts the bits a hardware implementation of spp16
needs (the simulator uses 64-bit C++ variables for convenience). Strides are
stored in 7 bits ($-64\ldots+63$ lines): a prefetch can never leave the 64-line
page, so larger strides can be represented by one saturated ``out of range''
code without losing a prefetch. Replacement uses true LRU ($\log_2$ of the
table size per entry).

\begin{table}[htbp]
\centering
\caption{Bit-level storage breakdown of the final SPP (\texttt{spp16}).}\label{tab:spp-storage}
\tablesetup
\begin{tabular}{llrrr}
\toprule
\rowcolor{tablehead}\colhead{Structure} & \colhead{Field} & \colhead{Bits/entry} & \colhead{Entries} & \colhead{Total bits} \\
\midrule
IP tracker     & Valid                          & 1  & 128 & 128 \\
               & IP tag (\texttt{MASK\_IP 0xFF}) & 8  & 128 & 1{,}024 \\
               & Last line address              & 58 & 128 & 7{,}424 \\
               & Stride history (3 $\times$ 7)  & 21 & 128 & 2{,}688 \\
               & Strides seen (0--3)            & 2  & 128 & 256 \\
               & Chooser                        & 2  & 128 & 256 \\
               & LRU                            & 7  & 128 & 896 \\
\midrule
Short table    & History (1 stride)             & 7  & 1{,}024 & 7{,}168 \\
(8 per IP)     & Predicted stride               & 7  & 1{,}024 & 7{,}168 \\
               & Valid                          & 1  & 1{,}024 & 1{,}024 \\
               & LRU                            & 3  & 1{,}024 & 3{,}072 \\
\midrule
Long table     & Signature (3 strides)          & 21 & 1{,}024 & 21{,}504 \\
(8 per IP)     & Predicted stride               & 7  & 1{,}024 & 7{,}168 \\
               & Valid                          & 1  & 1{,}024 & 1{,}024 \\
               & LRU                            & 3  & 1{,}024 & 3{,}072 \\
\midrule
Throttling     & Window fetched counter         & 9  & 1 & 9 \\
               & Window useful counter          & 10 & 1 & 10 \\
               & Current degree                 & 4  & 1 & 4 \\
\midrule
\multicolumn{4}{l}{\textbf{Total}} & \textbf{63{,}895 bits (7{,}987 B $\approx$ 7.8 KB)} \\
\bottomrule
\end{tabular}
\end{table}

\begin{table}[htbp]
\centering
\caption{Storage versus performance (same algorithm, different sizes).}\label{tab:storage-perf}
\tablesetup
\begin{tabular}{lrrrrr}
\toprule
\rowcolor{tablehead}\colhead{Version} & \colhead{Trackers} & \colhead{Short / long entries} & \colhead{Bits} & \colhead{Storage} & \colhead{Geomean} \\
\midrule
spp14 & 256 & 24 / 64 & 721{,}943 & 88.1\,KB & 1.0144 \\
spp15 & 128 & 8 / 16 & 98{,}711 & 12.0\,KB & 1.0129 \\
\textbf{spp16} & \textbf{128} & \textbf{8 / 8} & \textbf{63{,}895} & \textbf{7.8\,KB} & \textbf{1.0129} \\
\bottomrule
\end{tabular}
\end{table}

\paragraph{Further reduction.} Half of the remaining bits are the long tables
and the stored signatures. A single shared long table, direct-mapped by a
hash of the three-stride signature (e.g.\ 1024 entries of a 6-bit partial tag,
a 7-bit prediction and a valid bit, $\approx$14\,kbit), would replace the
per-IP long tables and their LRU bits and bring the design to roughly 5.5\,KB; IPs
with the same pattern would also share training. Storing the last address as a
page hash plus a 6-bit offset would save a further $\approx$5\,kbit.

\subsubsection{Results}\label{sec:task1-results}
\Cref{tab:task1-results} compares the baseline, the provided IP-stride
prefetcher and the final SPP; \cref{fig:t1-speedup,fig:t1-mpki,fig:t1-acc,fig:t1-cov}
plot the same data.

\begin{table}[htbp]
\centering
\caption{Task 1 results (L1D prefetching, LRU at L2). SPP = final design \texttt{spp16}.}\label{tab:task1-results}
\tablesetup
\begin{tabular}{llrrrrr}
\toprule
\rowcolor{tablehead}\colhead{Trace} & \colhead{Config} & \colhead{IPC} & \colhead{Speedup} & \colhead{L1D load MPKI} & \colhead{Accuracy (\%)} & \colhead{Coverage (\%)} \\
\midrule
                          & Baseline  & 0.2320 & 1.0000 & 10.68 & -- & -- \\
                         & IP-stride & 0.2330 & 1.0044 & 10.12 & 52.9 & 5.2 \\
\multirow{-3}{*}{Trace 1} & SPP       & 0.2339 & 1.0080 & 9.52 & 47.6 & 10.9 \\
\midrule
                          & Baseline  & 0.2506 & 1.0000 & 13.54 & -- & -- \\
                         & IP-stride & 0.2508 & 1.0007 & 13.44 & 38.0 & 0.7 \\
\multirow{-3}{*}{Trace 2} & SPP       & 0.2541 & 1.0141 & 12.91 & 51.6 & 4.6 \\
\midrule
                          & Baseline  & 0.3819 & 1.0000 & 3.40 & -- & -- \\
                         & IP-stride & 0.3827 & 1.0021 & 3.32 & 16.0 & 2.5 \\
\multirow{-3}{*}{Trace 3} & SPP       & 0.3882 & 1.0165 & 2.88 & 33.4 & 15.3 \\
\midrule
                             & IP-stride & -- & 1.0024 & 8.96 & 35.6 & 2.8 \\
\multirow{-2}{*}{Mean} & SPP       & -- & \textbf{1.0129} & \textbf{8.44} & \textbf{44.2} & \textbf{10.3} \\
\bottomrule
\end{tabular}
\end{table}

\resultfig{task1_speedup}{Speedup over the no-prefetch baseline (bars: IPC gain in \%; labels: speedup).}{fig:t1-speedup}
\resultfig{task1_l1d_mpki}{L1D load MPKI.}{fig:t1-mpki}
\resultfig{task1_accuracy}{Prefetch accuracy (useful / prefetches fetched from L2).}{fig:t1-acc}
\resultfig{task1_coverage}{L1D load-miss coverage.}{fig:t1-cov}

\subsubsection{Discussion}\label{sec:task1-discussion}
\paragraph{Why the absolute gains are small.} \Cref{tab:frontend} shows the
baseline bottlenecks. Traces~1 and~2 are dominated by instruction misses
(L1I MPKI 14.5 and 40.0, and more L2 instruction than data misses) and
Trace~1 also by branch mispredictions (16.1 MPKI); an L1D prefetcher cannot
touch these stalls. Trace~3 is the most data-bound, but its L1D misses only
3.4 times per kilo-instruction. Speedups of 1--2\% are therefore close to what
an L1D prefetcher can deliver here; relative to IP-stride the final SPP gains
5.4$\times$ more.

\begin{table}[htbp]
\centering
\caption{Baseline bottlenecks (no prefetching).}\label{tab:frontend}
\tablesetup\footnotesize
\begin{tabular}{lrrrrrr}
\toprule
\rowcolor{tablehead}\colhead{Trace} & \colhead{IPC} & \colhead{L1I MPKI} & \colhead{L1D load MPKI} & \colhead{L2 instr.\ MPKI} & \colhead{L2 data MPKI} & \colhead{Branch MPKI} \\
\midrule
Trace 1 & 0.232 & 14.5 & 10.7 & 10.4 & 7.3 & 16.1 \\
Trace 2 & 0.251 & 40.0 & 13.5 & 19.6 & 8.6 & 6.3 \\
Trace 3 & 0.382 & 2.1 & 3.4 & 1.1 & 2.8 & 7.2 \\
\bottomrule
\end{tabular}
\end{table}

\paragraph{Trace by trace.} \emph{Trace~1}: SPP doubles IP-stride's coverage
(10.9 vs.\ 5.2\%) at slightly lower accuracy; its predictable misses are mostly
regular streams, which the deeper walk and 128 trackers serve better.
\emph{Trace~2}: many live IPs; IP-stride's 24 trackers thrash (0.7\% coverage),
whereas SPP covers 4.6\% with higher accuracy (51.6\%), giving the largest
relative gain (about 19$\times$ IP-stride's speedup). \emph{Trace~3}: regular streams
whose prefetches were mostly late with IP-stride (53\% late); the deeper walk
cuts lateness to 18\%, and throttling raises accuracy from 16\% to 33\%,
giving the best result (1.0165, coverage 15.3\%).

\paragraph{Insights.} (i) Timeliness mattered more than prediction quality:
the largest single gains came from walking further ahead. (ii) For these
traces a short history covers most predictable misses; longer signatures are
more precise but cover less, so they must be a reversible fall-back, never a
one-way escalation. (iii) Throttling helps by \emph{reducing} the degree in bad
phases; beyond degree 8 extra lookahead is wasted inside a 4\,KB page.
(iv) Request redundancy remains: over 95\% of requests do not fetch a line
(2.2M requested vs.\ 97k fetched on Trace~1), mostly re-requests of lines
already covered by the previous walk; a per-IP
``already prefetched up to step $n$'' marker would remove them without changing
coverage.

\subsubsection{Full Comparison of All Versions}\label{sec:task1-full}
% make_table.py: layout=long (47 rows and 641pt wide -> longtable over pages; font \scriptsize to fit)
\FloatBarrier  % keep earlier floats before this inline table
\begingroup
\tablesetup \scriptsize
\begin{longtable}{llrrrrrrrrr}
\caption{Full comparison of all evaluated versions (spp13 omitted: identical to spp14 up to rounding). Coverage, accuracy and lateness are percentages; Fetched = prefetches that brought a line from L2; Requested = calls to \texttt{prefetch\_line}.}\label{tab:full}\\
\toprule
\rowcolor{tablehead}\colhead{Version} & \colhead{Trace} & \colhead{IPC} & \colhead{Speedup} & \colhead{L1D ld MPKI} & \colhead{Cov. (\%)} & \colhead{Acc. (\%)} & \colhead{Late (\%)} & \colhead{Fetched} & \colhead{Useful} & \colhead{Requested} \\
\midrule
\endfirsthead
\multicolumn{11}{l}{\small\itshape\color{muted}\tablename~\thetable{} (continued)}\\
\toprule
\rowcolor{tablehead}\colhead{Version} & \colhead{Trace} & \colhead{IPC} & \colhead{Speedup} & \colhead{L1D ld MPKI} & \colhead{Cov. (\%)} & \colhead{Acc. (\%)} & \colhead{Late (\%)} & \colhead{Fetched} & \colhead{Useful} & \colhead{Requested} \\
\midrule
\endhead
\midrule
\multicolumn{11}{r}{\small\itshape\color{muted}continued on next page}\\
\endfoot
\bottomrule
\endlastfoot
baseline & 1 & 0.2320 & 1.0000 & 10.68 & -- & -- & -- & 0 & 0 & 0 \\
baseline & 2 & 0.2506 & 1.0000 & 13.54 & -- & -- & -- & 0 & 0 & 0 \\
baseline & 3 & 0.3819 & 1.0000 & 3.40 & -- & -- & -- & 0 & 0 & 0 \\
ip\_stride & 1 & 0.2330 & 1.0044 & 10.12 & 5.2 & 52.9 & 14.7 & 30971 & 16378 & 298580 \\
ip\_stride & 2 & 0.2508 & 1.0007 & 13.44 & 0.7 & 38.0 & 6.0 & 11521 & 4380 & 152145 \\
ip\_stride & 3 & 0.3827 & 1.0021 & 3.32 & 2.5 & 16.0 & 52.7 & 17069 & 2731 & 200869 \\
spp1 & 1 & 0.2325 & 1.0020 & 10.43 & 2.3 & 60.7 & 27.2 & 10569 & 6412 & 76900 \\
spp2 & 1 & 0.2328 & 1.0033 & 10.24 & 4.1 & 55.4 & 14.1 & 22038 & 12213 & 289766 \\
spp2 & 2 & 0.2508 & 1.0011 & 13.46 & 0.6 & 57.2 & 5.2 & 4951 & 2834 & 89619 \\
spp2 & 3 & 0.3830 & 1.0029 & 3.29 & 3.3 & 22.1 & 36.8 & 14868 & 3293 & 180162 \\
spp3 & 1 & 0.2332 & 1.0052 & 9.92 & 7.1 & 48.5 & 11.5 & 46786 & 22670 & 616360 \\
spp3 & 2 & 0.2507 & 1.0007 & 13.43 & 0.8 & 32.9 & 5.0 & 16517 & 5438 & 249386 \\
spp3 & 3 & 0.3835 & 1.0043 & 3.19 & 6.2 & 18.3 & 35.3 & 35011 & 6412 & 548178 \\
spp4 & 1 & 0.2331 & 1.0049 & 10.01 & 6.3 & 50.3 & 11.1 & 39861 & 20060 & 507874 \\
spp4 & 2 & 0.2508 & 1.0007 & 13.43 & 0.8 & 33.5 & 5.0 & 15812 & 5301 & 239275 \\
spp4 & 3 & 0.3832 & 1.0035 & 3.26 & 4.2 & 19.5 & 31.7 & 22793 & 4453 & 437681 \\
spp5 & 2 & 0.2518 & 1.0049 & 13.28 & 1.9 & 43.9 & 2.6 & 21355 & 9381 & 537738 \\
spp6 & 1 & 0.2322 & 1.0007 & 10.62 & 0.6 & 78.2 & 7.8 & 2460 & 1924 & 38650 \\
spp6 & 2 & 0.2512 & 1.0026 & 13.48 & 0.4 & 75.9 & 5.1 & 2419 & 1837 & 112777 \\
spp6 & 3 & 0.3819 & 1.0001 & 3.40 & 0.1 & 60.9 & 18.8 & 184 & 112 & 33219 \\
spp7 & 1 & 0.2334 & 1.0058 & 9.70 & 9.2 & 45.6 & 5.4 & 77930 & 35559 & 1522111 \\
spp7 & 2 & 0.2521 & 1.0060 & 13.22 & 2.4 & 34.8 & 1.5 & 36327 & 12641 & 935257 \\
spp7 & 3 & 0.3857 & 1.0100 & 3.08 & 9.5 & 23.2 & 11.0 & 44226 & 10271 & 964729 \\
spp8 & 1 & 0.2336 & 1.0070 & 9.41 & 11.9 & 43.5 & 4.4 & 125677 & 54630 & 3126858 \\
spp8 & 2 & 0.2543 & 1.0147 & 12.88 & 4.9 & 42.2 & 1.4 & 63014 & 26612 & 2049660 \\
spp8 & 3 & 0.3887 & 1.0179 & 2.82 & 17.1 & 21.8 & 14.5 & 83769 & 18243 & 2403452 \\
spp9 & 1 & 0.2334 & 1.0060 & 9.81 & 8.1 & 56.5 & 5.5 & 55727 & 31500 & 1748725 \\
spp9 & 2 & 0.2545 & 1.0155 & 12.85 & 5.0 & 71.3 & 1.7 & 29353 & 20942 & 1371169 \\
spp9 & 3 & 0.3855 & 1.0095 & 3.07 & 9.7 & 27.1 & 18.4 & 34760 & 9423 & 1709591 \\
spp10 & 1 & 0.2337 & 1.0073 & 9.37 & 12.2 & 44.5 & 4.4 & 124230 & 55233 & 3065988 \\
spp10 & 2 & 0.2543 & 1.0148 & 12.85 & 5.1 & 45.3 & 1.4 & 64157 & 29037 & 2057559 \\
spp10 & 3 & 0.3889 & 1.0183 & 2.80 & 17.6 & 24.0 & 14.8 & 75690 & 18144 & 2223233 \\
spp11 & 1 & 0.2337 & 1.0073 & 9.40 & 12.0 & 45.0 & 4.5 & 119368 & 53732 & 2952185 \\
spp11 & 2 & 0.2543 & 1.0148 & 12.85 & 5.0 & 47.2 & 1.6 & 58896 & 27819 & 1951350 \\
spp11 & 3 & 0.3887 & 1.0177 & 2.84 & 16.6 & 23.7 & 13.9 & 71822 & 17041 & 2119438 \\
spp12 & 1 & 0.2335 & 1.0065 & 9.38 & 12.2 & 43.1 & 4.3 & 129922 & 56054 & 3504498 \\
spp12 & 2 & 0.2543 & 1.0149 & 12.84 & 5.1 & 45.2 & 1.4 & 64064 & 28950 & 2492034 \\
spp12 & 3 & 0.3897 & 1.0204 & 2.77 & 18.6 & 33.5 & 15.3 & 53851 & 18063 & 1480879 \\
spp14 & 1 & 0.2339 & 1.0079 & 9.39 & 12.0 & 46.6 & 4.9 & 114478 & 53306 & 2748438 \\
spp14 & 2 & 0.2545 & 1.0157 & 12.82 & 5.3 & 50.2 & 1.6 & 56174 & 28206 & 1829211 \\
spp14 & 3 & 0.3894 & 1.0196 & 2.80 & 17.8 & 35.2 & 16.5 & 48678 & 17123 & 1334309 \\
spp15 & 1 & 0.2339 & 1.0081 & 9.52 & 10.9 & 47.5 & 5.4 & 97277 & 46244 & 2206887 \\
spp15 & 2 & 0.2541 & 1.0141 & 12.91 & 4.7 & 52.2 & 2.1 & 38336 & 19996 & 1169160 \\
spp15 & 3 & 0.3882 & 1.0165 & 2.89 & 15.2 & 33.6 & 18.3 & 43637 & 14662 & 957902 \\
spp16 & 1 & 0.2339 & 1.0080 & 9.52 & 10.9 & 47.6 & 5.2 & 97421 & 46334 & 2200401 \\
spp16 & 2 & 0.2541 & 1.0141 & 12.91 & 4.6 & 51.6 & 2.0 & 38673 & 19964 & 1194763 \\
spp16 & 3 & 0.3882 & 1.0165 & 2.88 & 15.3 & 33.4 & 17.8 & 44075 & 14707 & 907982 \\
\end{longtable}
\endgroup




\end{document}
